\chapter{Number Theory: Primes, GCD, and Euclid}

Almost every public-key system in use today --- RSA, Diffie-Hellman, elliptic curves, the
signatures that protect your bank login --- rests on a small pile of facts about the integers
that were already old when Euclid wrote them down around 300 BCE. The integers look like the
most boring objects in mathematics: you can add them, multiply them, and that is about it. But
hidden inside that plain arithmetic is a deep asymmetry. Multiplying two large primes is fast and
easy; tearing the product back apart into its factors is, as far as anyone knows, ruinously slow.
Modern cryptography lives inside that gap.

This chapter builds the toolkit. We start with divisibility and primes, prove that every integer
factors into primes in exactly one way, and meet the Euclidean algorithm --- a procedure so
efficient and so old that it is genuinely startling. Extended Euclid hands us modular inverses,
which is how RSA decryption keys get computed. Then come the two theorems that make exponentiation
modulo a number behave predictably: Fermat's little theorem and its generalization, Euler's
theorem, both proved from scratch. The Euler totient counts the units; the Chinese Remainder
Theorem lets us split a hard modular problem into easy independent pieces (and speeds up RSA by a
factor of four); primitive roots give us the cyclic structure that Diffie-Hellman needs. Everything
is done with tiny numbers you can check on paper, because tracing real little numbers is the only
way the abstractions ever become yours.

\chapterline

\section{Divisibility}
\label{ch04:divisibility}

We work inside $\Z = \set{\dots, -2, -1, 0, 1, 2, \dots}$, the integers.

\begin{definition}[Divisibility]
\label{ch04:def-divides}
For integers $a$ and $b$, we say $a$ \emph{divides} $b$, written $a \mid b$, if there exists an
integer $k$ with $b = ak$. We call $a$ a \emph{divisor} (or \emph{factor}) of $b$, and $b$ a
\emph{multiple} of $a$. If no such $k$ exists we write $a \nmid b$.
\end{definition}

So $3 \mid 12$ because $12 = 3 \cdot 4$, and $7 \nmid 12$ because no integer multiple of $7$ lands
on $12$. Notice $a \mid 0$ for every $a$ (take $k = 0$), and $0 \mid b$ only when $b = 0$. The
relation $a \mid b$ says nothing about size beyond this: $-3 \mid 12$ just as much as $3$ does.

A few facts fall straight out of the definition, and we will lean on them constantly.

\begin{proposition}[Basic divisibility rules]
\label{ch04:prop-div-rules}
Let $a, b, c$ be integers.
\begin{enumerate}
\item If $a \mid b$ and $b \mid c$ then $a \mid c$ (transitivity).
\item If $a \mid b$ and $a \mid c$ then $a \mid (bx + cy)$ for all integers $x, y$.
\item If $a \mid b$ and $b \neq 0$ then $\abs{a} \le \abs{b}$.
\end{enumerate}
\end{proposition}

\begin{proof}
(1) Write $b = ak$ and $c = bm$. Then $c = a(km)$, so $a \mid c$.
(2) Write $b = ak$ and $c = a\ell$. Then $bx + cy = a(kx + \ell y)$, an integer multiple of $a$.
(3) Write $b = ak$ with $k \neq 0$ (since $b \neq 0$). Then $\abs{b} = \abs{a}\abs{k} \ge \abs{a}$
because $\abs{k} \ge 1$.
\end{proof}

Part (2) is the workhorse. The set of all integer combinations $bx + cy$ is exactly the set of
multiples of the greatest common divisor of $b$ and $c$, a fact we will prove and then use over and
over.

\subsection*{Division with remainder}

You cannot always divide one integer evenly by another, but you can always divide \emph{with a
remainder}, and that remainder is unique. This is the single most important computational fact in
the chapter.

\begin{theorem}[Division algorithm]
\label{ch04:thm-division}
For integers $a$ and $b$ with $b > 0$, there exist \emph{unique} integers $q$ (the quotient) and
$r$ (the remainder) such that
\[
  a = bq + r, \qquad 0 \le r < b.
\]
\end{theorem}

\begin{proof}
\emph{Existence.} Consider the set $S = \set{a - bq : q \in \Z,\ a - bq \ge 0}$. It is nonempty:
taking $q = -\abs{a}$ gives $a - bq = a + b\abs{a} \ge a + \abs{a} \ge 0$. By the well-ordering
principle, $S$ has a least element $r = a - bq$ for some $q$. We have $r \ge 0$ by construction.
If $r \ge b$, then $r - b = a - b(q+1) \ge 0$ would be a smaller element of $S$, contradicting
minimality. So $0 \le r < b$.

\emph{Uniqueness.} Suppose $a = bq_1 + r_1 = bq_2 + r_2$ with both remainders in $[0, b)$.
Subtracting, $b(q_1 - q_2) = r_2 - r_1$, so $b \mid (r_2 - r_1)$. But $\abs{r_2 - r_1} < b$, and the
only multiple of $b$ with absolute value below $b$ is $0$. Hence $r_1 = r_2$, and then
$q_1 = q_2$.
\end{proof}

This $r$ is what your programming language writes as \texttt{a \% b} (when $b > 0$ and the language
follows the mathematical convention). It is the bridge to modular arithmetic: $a \equiv r
\pmod{b}$.

\begin{intuitionbox}
Lay out a number line and mark the multiples of $b$: $\dots, -b, 0, b, 2b, \dots$. Every integer
$a$ sits in exactly one of the half-open gaps $[bq, bq+b)$. The remainder $r$ is how far $a$ is
past the multiple of $b$ on its left. There is one gap per integer, so $q$ and $r$ are forced ---
that is the uniqueness.
\end{intuitionbox}

\section{Primes and the Fundamental Theorem}
\label{ch04:primes}

\begin{definition}[Prime and composite]
\label{ch04:def-prime}
An integer $p > 1$ is \emph{prime} if its only positive divisors are $1$ and $p$. An integer
$n > 1$ that is not prime is \emph{composite}: it factors as $n = ab$ with $1 < a, b < n$. By
convention $1$ is neither prime nor composite.
\end{definition}

The first primes are $2, 3, 5, 7, 11, 13, 17, 19, 23, 29, \dots$. Two is the only even prime, which
makes it the perennial special case in proofs and code.

\begin{historybox}
Euclid's \emph{Elements}, Book IX, Proposition 20, states: ``Prime numbers are more than any
assigned multitude of prime numbers.'' That is, there are infinitely many primes. His argument is
the one every mathematician learns first, and it has not been improved in 2300 years. Suppose you
had a finite list $p_1, \dots, p_k$ of all the primes. Form $N = p_1 p_2 \cdots p_k + 1$. Dividing
$N$ by any $p_i$ leaves remainder $1$, so no $p_i$ divides $N$. But $N > 1$ has \emph{some} prime
factor (we prove this below), and that factor is not on the list --- contradiction. The list could
never have been complete.
\end{historybox}

\subsection*{Every integer has a prime factor}

\begin{lemma}
\label{ch04:lem-has-prime-factor}
Every integer $n > 1$ has at least one prime divisor.
\end{lemma}

\begin{proof}
Let $d$ be the smallest divisor of $n$ that is greater than $1$ (such divisors exist because $n$
itself is one, and the well-ordering principle gives a least one). If $d$ were composite, $d = ab$
with $1 < a < d$, then $a$ would be a divisor of $n$ smaller than $d$ but still above $1$,
contradicting minimality. So $d$ is prime.
\end{proof}

\subsection*{Euclid's lemma}

The keystone result --- the one that makes factorization \emph{unique} --- is deceptively simple to
state and a little subtle to prove. It says primes cannot be split across a product.

\begin{lemma}[Euclid's lemma]
\label{ch04:lem-euclid}
If $p$ is prime and $p \mid ab$, then $p \mid a$ or $p \mid b$.
\end{lemma}

We will prove it cleanly once we have Bezout's identity from the Euclidean algorithm
(Section~\ref{ch04:euclid}), because the natural proof uses the fact that $\gcd(p, a) = 1$ whenever
$p \nmid a$. For now, accept it and watch what it buys us. Note the primality is essential:
$6 \mid 4 \cdot 9 = 36$, yet $6 \nmid 4$ and $6 \nmid 9$. A composite number can hide inside a
product by lending one factor to each side.

\begin{theorem}[Fundamental Theorem of Arithmetic]
\label{ch04:thm-fta}
Every integer $n > 1$ can be written as a product of primes,
\[
  n = p_1^{e_1} p_2^{e_2} \cdots p_k^{e_k},
\]
and this factorization is unique up to the order of the factors.
\end{theorem}

\begin{proof}
\emph{Existence} is strong induction. If $n$ is prime, it is its own factorization. Otherwise
$n = ab$ with $1 < a, b < n$; by induction both $a$ and $b$ are products of primes, and
concatenating those products factors $n$.

\emph{Uniqueness.} Suppose $n$ has two prime factorizations,
\[
  p_1 p_2 \cdots p_s = q_1 q_2 \cdots q_t,
\]
listing primes with repetition. The prime $p_1$ divides the left side, hence the right side, so by
Euclid's lemma (applied repeatedly) $p_1 \mid q_j$ for some $j$. Since $q_j$ is prime its only
divisors are $1$ and $q_j$, forcing $p_1 = q_j$. Cancel that common prime from both sides and
repeat. Each step removes one prime from each list; the lists must run out together, so $s = t$ and
the multisets of primes coincide.
\end{proof}

\begin{focusbox}
Uniqueness is not obvious and it is not free --- it is exactly Euclid's lemma in disguise. There are
number systems that look like the integers but where factorization is \emph{not} unique. In the
ring $\Z[\sqrt{-5}]$ the number $6$ factors as $2 \cdot 3$ and also as $(1+\sqrt{-5})(1-\sqrt{-5})$,
with all four factors ``prime'' in the sense of having no smaller factorization. Euclid's lemma
fails there. The integers are special, and cryptography quietly depends on that specialness every
time it assumes an RSA modulus has exactly one factorization $n = pq$.
\end{focusbox}

\section{Greatest Common Divisor and the Euclidean Algorithm}
\label{ch04:euclid}

\begin{definition}[GCD]
\label{ch04:def-gcd}
The \emph{greatest common divisor} $\gcd(a, b)$ of two integers, not both zero, is the largest
positive integer that divides both. When $\gcd(a, b) = 1$ we call $a$ and $b$ \emph{coprime} (or
relatively prime): they share no prime factor.
\end{definition}

You could compute $\gcd(a,b)$ by factoring both numbers and multiplying the shared prime powers ---
$\gcd(60, 84) = \gcd(2^2 \cdot 3 \cdot 5,\ 2^2 \cdot 3 \cdot 7) = 2^2 \cdot 3 = 12$. But factoring is
slow. Euclid found a way to get the gcd \emph{without} factoring at all, and it is fast.

\subsection*{The algorithm}

The whole thing rests on one observation.

\begin{lemma}
\label{ch04:lem-gcd-mod}
For $b > 0$, $\gcd(a, b) = \gcd(b, a \bmod b)$.
\end{lemma}

\begin{proof}
Write $a = bq + r$ with $r = a \bmod b$. Any common divisor of $a$ and $b$ divides
$a - bq = r$ (by Proposition~\ref{ch04:prop-div-rules}, part 2), so it is a common divisor of $b$
and $r$. Conversely any common divisor of $b$ and $r$ divides $bq + r = a$, so it is a common
divisor of $a$ and $b$. The two pairs have identical sets of common divisors, hence the same
greatest one.
\end{proof}

So replace $(a, b)$ by $(b, a \bmod b)$ and repeat. The second coordinate strictly shrinks and stays
nonnegative, so it must hit $0$. When it does, $\gcd(g, 0) = g$, and that $g$ is the answer.

\begin{example}[$\gcd(1071, 462)$ by hand]
\label{ch04:ex-gcd-1071}
\[
\begin{aligned}
1071 &= 2 \cdot 462 + 147 \\
 462 &= 3 \cdot 147 + 21 \\
 147 &= 7 \cdot 21 + 0.
\end{aligned}
\]
The last nonzero remainder is $21$, so $\gcd(1071, 462) = 21$. Three divisions. Factoring
$1071 = 3 \cdot 357 = 3 \cdot 3 \cdot 7 \cdot 17$ would have taken longer and given the same
answer.
\end{example}

\begin{figure}[H]
\centering
\begin{tikzpicture}[scale=0.18,font=\small]
  % 48 x 18 rectangle outline
  \draw[thick,black] (0,0) rectangle (48,18);

  % First division: 48 = 2 * 18 + 12
  \draw[thick,sourceblue!80,fill=sourceblue!15] (0,0) rectangle (18,18);
  \node[sourceblue!90!black,font=\footnotesize\bfseries] at (9,9) {Square $18 \times 18$};
  
  \draw[thick,sourceblue!80,fill=sourceblue!15] (18,0) rectangle (36,18);
  \node[sourceblue!90!black,font=\footnotesize\bfseries] at (27,9) {Square $18 \times 18$};

  % Second division: 18 = 1 * 12 + 6
  \draw[thick,targetgreen!80,fill=targetgreen!20] (36,0) rectangle (48,12);
  \node[targetgreen!90!black,font=\scriptsize\bfseries] at (42,6) {$12 \times 12$};

  % Third division: 12 = 2 * 6 + 0
  \draw[thick,quotientorange!80,fill=quotientorange!25] (36,12) rectangle (42,18);
  \node[quotientorange!90!black,font=\tiny\bfseries] at (39,15) {$6 \times 6$};

  \draw[thick,quotientorange!80,fill=quotientorange!25] (42,12) rectangle (48,18);
  \node[quotientorange!90!black,font=\tiny\bfseries] at (45,15) {$6 \times 6$};

  % Dimension arrows
  \draw[<->,thick,black] (0,-2.5) -- node[below,font=\footnotesize\bfseries]{Length $a = 48$} (48,-2.5);
  \draw[<->,thick,black] (-2.5,0) -- node[left,font=\footnotesize\bfseries]{Width $b = 18$} (0,18);

  % Annotations
  \node[right,font=\footnotesize,align=left] at (51,9) {
    \textbf{Geometric Euclidean Algorithm:}\\
    $\bullet$ \textbf{Step 1:} Cut two $18\times 18$ squares $\implies$ Remainder $12\times 18$\\
    $\bullet$ \textbf{Step 2:} Cut one $12\times 12$ square $\implies$ Remainder $6\times 12$\\
    $\bullet$ \textbf{Step 3:} Cut two $6\times 6$ squares $\implies$ Zero remainder!\\
    \textbf{Conclusion:} Largest square tile is $6\times 6 \implies \gcd(48,18) = 6$.
  };
\end{tikzpicture}
\caption{The Euclidean algorithm as geometric rectangle tiling for $\gcd(48, 18) = 6$. At each step, the largest possible square is cut out from the remaining rectangle. The smallest square that tiles the final remainder without gaps perfectly tiles all preceding squares, revealing the greatest common divisor.}
\label{ch04:fig-euclid-tiling}
\end{figure}

\begin{connectionbox}
How fast is this? The remainders shrink quickly --- in fact each two steps at least halve the
smaller number, so the number of divisions is $O(\log b)$. The slowest possible inputs are
\emph{consecutive Fibonacci numbers}: $\gcd(F_{n+1}, F_n)$ takes exactly $n-1$ divisions, because
each quotient is forced to be $1$. This is Lame's theorem (1844), historically the first time
anyone bounded an algorithm's running time in terms of input size. The companion algorithms book
works out the $O(\log b)$ bound and the Fibonacci worst case in detail; the upshot for us is that
gcd of two thousand-bit numbers finishes in a few thousand machine operations, instantly.
\end{connectionbox}

\subsection*{Bezout's identity and the extended algorithm}

The Euclidean algorithm gives more than the gcd. By unwinding the divisions we can write the gcd as
an integer combination of the inputs.

\begin{theorem}[Bezout's identity]
\label{ch04:thm-bezout}
For integers $a, b$ not both zero, there exist integers $x, y$ with
\[
  ax + by = \gcd(a, b).
\]
Moreover $\gcd(a,b)$ is the \emph{smallest} positive value taken by $ax + by$ over all integers
$x, y$.
\end{theorem}

\begin{proof}
Let $d$ be the smallest positive integer of the form $ax + by$ (the set of positive such values is
nonempty --- it contains $\abs{a}$ when $a \neq 0$ --- so well-ordering gives a least element). We
show $d = \gcd(a, b)$.

First, $d$ divides both $a$ and $b$. Divide $a = dq + r$ with $0 \le r < d$. Then
\[
  r = a - dq = a - q(ax + by) = a(1 - qx) + b(-qy),
\]
which is again of the form $a x' + b y'$. Since $0 \le r < d$ and $d$ is the smallest
\emph{positive} such value, we must have $r = 0$, i.e. $d \mid a$. The same argument gives
$d \mid b$. So $d$ is a common divisor.

Second, $d$ is the greatest. Any common divisor $c$ of $a$ and $b$ divides $ax + by = d$, so
$c \le d$. Hence $d = \gcd(a, b)$.
\end{proof}

The coefficients $x, y$ are called \emph{Bezout coefficients}, and we compute them by running
Euclid forward while carrying the combinations along. That is the \emph{extended Euclidean
algorithm}. Continuing Example~\ref{ch04:ex-gcd-1071}, we back-substitute:
\[
\begin{aligned}
21 &= 462 - 3 \cdot 147 \\
   &= 462 - 3 \cdot (1071 - 2 \cdot 462) \\
   &= 7 \cdot 462 - 3 \cdot 1071.
\end{aligned}
\]
Check: $7 \cdot 462 - 3 \cdot 1071 = 3234 - 3213 = 21$. So $x = -3$, $y = 7$.

\begin{figure}[H]
\centering
\begin{tikzpicture}[font=\small,>=stealth,scale=0.95]
  % Forward Descent Card (Euclid)
  \begin{scope}[shift={(0,0)}]
    \draw[thick,sourceblue!50,fill=sourceblue!5,rounded corners=6pt] (-0.3,-2.8) rectangle (5.2,1.8);
    \node[font=\bfseries\footnotesize,sourceblue!90!black] at (2.45,1.4) {1. Forward Division Descent};

    \node[draw=sourceblue!60,fill=white,rounded corners=2pt,thick,minimum width=4.8cm,minimum height=0.7cm,font=\footnotesize] (eq1) at (2.45,0.6) {
      $1071 = 2 \cdot 462 + \mathbf{147}$ \hfill \tiny $q_1=2$
    };
    \node[draw=sourceblue!60,fill=white,rounded corners=2pt,thick,minimum width=4.8cm,minimum height=0.7cm,font=\footnotesize] (eq2) at (2.45,-0.3) {
      $462 = 3 \cdot 147 + \mathbf{21}$ \hfill \tiny $q_2=3$
    };
    \node[draw=targetgreen!70,fill=targetgreen!10,rounded corners=2pt,thick,minimum width=4.8cm,minimum height=0.7cm,font=\footnotesize\bfseries] (eq3) at (2.45,-1.2) {
      $147 = 7 \cdot \mathbf{21} + 0$ \hfill \tiny $\gcd = 21$
    };

    \draw[->,very thick,sourceblue] (2.45,-1.7) -- node[right,font=\tiny\bfseries]{$\gcd(1071,462) = 21$} (2.45,-2.4);
  \end{scope}

  % Backward Substitution Card (Bézout)
  \begin{scope}[shift={(6.2,0)}]
    \draw[thick,quotientorange!50,fill=quotientorange!5,rounded corners=6pt] (-0.3,-2.8) rectangle (5.6,1.8);
    \node[font=\bfseries\footnotesize,quotientorange!90!black] at (2.65,1.4) {2. Backward Substitution Ascent};

    \node[draw=quotientorange!60,fill=white,rounded corners=2pt,thick,minimum width=5.2cm,minimum height=0.7cm,font=\scriptsize] (b1) at (2.65,0.6) {
      $21 = 462 - 3 \cdot \mathbf{147}$
    };
    \node[draw=quotientorange!60,fill=white,rounded corners=2pt,thick,minimum width=5.2cm,minimum height=0.7cm,font=\scriptsize] (b2) at (2.65,-0.3) {
      $21 = 462 - 3 \cdot (1071 - 2 \cdot 462)$
    };
    \node[draw=structurepurple!70,fill=structurepurple!10,rounded corners=2pt,thick,minimum width=5.2cm,minimum height=0.7cm,font=\scriptsize\bfseries] (b3) at (2.65,-1.2) {
      $21 = \mathbf{7} \cdot 462 - \mathbf{3} \cdot 1071$
    };

    \node[draw=structurepurple!80,fill=white,rounded corners=2pt,thick,font=\tiny\bfseries,align=center] at (2.65,-2.1) {
      B\'ezout: $a(-3) + b(7) = 21$\\
      $x = -3,\quad y = 7$
    };
  \end{scope}

  % Connection between two columns
  \draw[->,very thick,dashed,quotientorange] (5.0,-1.2) to[bend left=25] (6.0,0.6);
\end{tikzpicture}
\caption{The Extended Euclidean algorithm ladder for $\gcd(1071, 462) = 21$. Left: Forward divisions isolate the remainder at each step until reaching the greatest common divisor 21. Right: Backward substitution unwinds the ladder, replacing intermediate remainders until 21 is expressed directly in terms of the initial inputs $1071(-3) + 462(7) = 21$.}
\label{ch04:fig-bezout-ladder}
\end{figure}

\begin{whybox}
Now Euclid's lemma (Lemma~\ref{ch04:lem-euclid}) drops out for free. Suppose $p$ is prime,
$p \mid ab$, and $p \nmid a$. Then $\gcd(p, a) = 1$ --- the only divisors of $p$ are $1$ and $p$,
and $p$ is ruled out --- so Bezout gives integers $x, y$ with $px + ay = 1$. Multiply by $b$:
\[
  pbx + aby = b.
\]
Now $p \mid pbx$ directly (since $p$ is a factor), and $p \mid aby$ because $p \mid ab$. So $p$ divides the left side,
hence $p \mid b$. That is exactly the claim: if a prime divides a product and misses one factor, it
must hit the other. The whole edifice of unique factorization rests on this two-line argument.
\end{whybox}

\subsection*{Computing modular inverses}

Here is where extended Euclid earns its keep in cryptography. In modular arithmetic the analogue of
``dividing by $a$'' is multiplying by the \emph{inverse} of $a$ modulo $n$: an integer $a^{-1}$ with
$a \cdot a^{-1} \equiv 1 \pmod n$.

\begin{theorem}[Existence of inverses]
\label{ch04:thm-inverse}
An integer $a$ has a multiplicative inverse modulo $n$ if and only if $\gcd(a, n) = 1$. When it
exists the inverse is unique modulo $n$.
\end{theorem}

\begin{proof}
If $\gcd(a, n) = 1$, Bezout gives $ax + ny = 1$, so $ax \equiv 1 \pmod n$ and $x$ is the inverse.
Conversely if $ax \equiv 1 \pmod n$ then $ax - 1 = ny$ for some $y$, i.e. $ax - ny = 1$; any common
divisor of $a$ and $n$ divides the left side and so divides $1$, forcing $\gcd(a, n) = 1$.

Uniqueness: if $ax \equiv 1$ and $ax' \equiv 1$, multiply the first by $x'$ to get
$x' \equiv x'ax \equiv x \pmod n$.
\end{proof}

So computing an inverse is exactly running extended Euclid and reading off $x$.

\begin{example}[Inverse of $17$ modulo $43$]
\label{ch04:ex-inverse}
Run Euclid on $(43, 17)$:
\[
\begin{aligned}
43 &= 2 \cdot 17 + 9 \\
17 &= 1 \cdot 9 + 8 \\
 9 &= 1 \cdot 8 + 1 \\
 8 &= 8 \cdot 1 + 0.
\end{aligned}
\]
The gcd is $1$, so the inverse exists. Back-substitute:
\[
\begin{aligned}
1 &= 9 - 1 \cdot 8 \\
  &= 9 - (17 - 9) = 2 \cdot 9 - 17 \\
  &= 2(43 - 2 \cdot 17) - 17 = 2 \cdot 43 - 5 \cdot 17.
\end{aligned}
\]
So $-5 \cdot 17 \equiv 1 \pmod{43}$, meaning $17^{-1} \equiv -5 \equiv 38 \pmod{43}$.
Check: $17 \cdot 38 = 646 = 15 \cdot 43 + 1$. Correct.
\end{example}

\begin{implementbox}
Extended Euclid in Python. The iterative form below tracks the Bezout coefficients in lockstep with
the remainders, so it never recurses and never overflows for big integers.

\begin{lstlisting}[language=Python]
def egcd(a, b):
    # returns (g, x, y) with a*x + b*y = g = gcd(a, b)
    old_r, r = a, b
    old_s, s = 1, 0
    old_t, t = 0, 1
    while r != 0:
        q = old_r // r
        old_r, r = r, old_r - q * r
        old_s, s = s, old_s - q * s
        old_t, t = t, old_t - q * t
    return old_r, old_s, old_t      # g, x, y

def inverse(a, n):
    g, x, _ = egcd(a % n, n)
    if g != 1:
        raise ValueError("no inverse: gcd(a, n) != 1")
    return x % n                    # normalize into [0, n)

print(inverse(17, 43))              # -> 38
\end{lstlisting}

\textbf{Learning only.} This is a clear, correct reference --- but it is not constant-time. The
number of loop iterations and the sizes of the intermediates depend on the secret input, which
leaks timing. Production libraries (for instance when computing an RSA private exponent) use
constant-time inversion or blinding. Use your platform's vetted big-integer routines, never a
hand-rolled loop, for real keys.
\end{implementbox}

\begin{applicationbox}
RSA key generation \emph{is} an extended-Euclid computation. You pick a public exponent $e$
(commonly $65537$) and must produce the private exponent $d$ with $ed \equiv 1 \pmod{\lambda(n)}$,
where $\lambda(n)$ is computed from the secret primes. That $d$ comes straight out of the extended
Euclidean algorithm applied to $e$ and $\lambda(n)$. If $\gcd(e, \lambda(n)) \neq 1$ the inverse
fails to exist and the key generator must pick different primes --- which is exactly why $e = 65537$
is convenient: it is prime, so it is coprime to $\lambda(n)$ unless $\lambda(n)$ happens to be a
multiple of $65537$, a vanishingly rare event you simply reject and retry.
\end{applicationbox}

\section{Congruences and the Units of $\Zn{n}$}
\label{ch04:congruences}

Recall $a \equiv b \pmod n$ means $n \mid (a - b)$. This relation partitions $\Z$ into $n$
\emph{residue classes}, and arithmetic respects them: if $a \equiv a'$ and $b \equiv b'$ then
$a + b \equiv a' + b'$ and $ab \equiv a'b' \pmod n$. The set of classes is written $\Zn{n}$, and we
usually name them by their least nonnegative representatives $\set{0, 1, \dots, n-1}$.

Not every class has a multiplicative inverse. By Theorem~\ref{ch04:thm-inverse}, the class of $a$ is
invertible exactly when $\gcd(a, n) = 1$. These invertible classes are the \emph{units}, and they
form a group under multiplication:
\[
  \Zns{n} = \set{a \in \set{1, \dots, n-1} : \gcd(a, n) = 1}.
\]

\begin{example}[Units modulo $12$]
\label{ch04:ex-units-12}
The classes coprime to $12 = 2^2 \cdot 3$ are $\set{1, 5, 7, 11}$. So $\abs{\Zns{12}} = 4$. Every
one of these is its own inverse: $5 \cdot 5 = 25 \equiv 1$, $7 \cdot 7 = 49 \equiv 1$,
$11 \cdot 11 = 121 \equiv 1 \pmod{12}$. The non-units $\set{2,3,4,6,8,9,10}$ have no inverse; for
instance $2 x \equiv 1 \pmod{12}$ is unsolvable because $2x$ is always even.
\end{example}

\subsection*{Euler's totient}

\begin{definition}[Euler totient]
\label{ch04:def-totient}
$\varphi(n)$ is the number of integers in $\set{1, \dots, n}$ that are coprime to $n$. Equivalently
$\varphi(n) = \abs{\Zns{n}}$, the size of the unit group.
\end{definition}

For a prime $p$, every one of $1, 2, \dots, p-1$ is coprime to $p$, so $\varphi(p) = p - 1$. For a
prime power $p^k$, the only integers in $\set{1, \dots, p^k}$ sharing a factor with $p^k$ are the
multiples of $p$, of which there are $p^{k-1}$. Hence
\[
  \varphi(p^k) = p^k - p^{k-1} = p^{k-1}(p - 1).
\]

The totient is \emph{multiplicative}: if $\gcd(m, n) = 1$ then $\varphi(mn) = \varphi(m)\varphi(n)$.
We prove this cleanly with the Chinese Remainder Theorem in Section~\ref{ch04:crt}. Combining the
two facts gives the general formula from the prime factorization
$n = p_1^{e_1} \cdots p_k^{e_k}$:
\[
  \varphi(n) = \prod_{i=1}^{k} p_i^{e_i - 1}(p_i - 1)
             = n \prod_{p \mid n}\left(1 - \frac{1}{p}\right).
\]

\begin{example}[$\varphi$ of an RSA-style modulus]
\label{ch04:ex-totient-rsa}
Take the textbook toy $n = 61 \cdot 53 = 3233$ with $p = 61$, $q = 53$ both prime. Then
\[
  \varphi(3233) = (61 - 1)(53 - 1) = 60 \cdot 52 = 3120.
\]
Knowing the two factors makes $\varphi(n)$ a one-line computation. Without them, computing
$\varphi(n)$ is as hard as factoring $n$ --- and that hardness is what keeps the RSA private key
private. Anyone who can find $\varphi(3233) = 3120$ can recover the prime factors: from
$p + q = n - \varphi(n) + 1 = 3233 - 3120 + 1 = 114$ and $pq = 3233$, solve the quadratic
$t^2 - 114 t + 3233 = 0$ to get $t = 61, 53$.
\end{example}

\begin{trapbox}
The formula $\varphi(pq) = (p-1)(q-1)$ requires $p$ and $q$ to be \emph{distinct} primes. If a buggy
key generator accidentally picks $p = q$, then $n = p^2$ and $\varphi(n) = p(p-1)$, not $(p-1)^2$;
worse, $n$ is now a perfect square and immediately factored by taking its integer square root. Real
generators explicitly check $p \neq q$. The same caution applies to the modern RSA convention of
using $\lambda(n) = \lcm(p-1, q-1)$ in place of $\varphi(n)$ --- it is a divisor of $\varphi(n)$, it
still works, and it yields a smaller private exponent.
\end{trapbox}

\section{Fermat's Little Theorem and Euler's Theorem}
\label{ch04:fermat-euler}

We now reach the two results that make modular exponentiation predictable --- the foundation of
RSA and of every primality test worth knowing.

\begin{theorem}[Fermat's little theorem]
\label{ch04:thm-flt}
If $p$ is prime and $p \nmid a$, then
\[
  a^{p-1} \equiv 1 \pmod p.
\]
Equivalently, for \emph{every} integer $a$, $a^p \equiv a \pmod p$.
\end{theorem}

\begin{proof}
Consider the $p - 1$ nonzero residues $1, 2, \dots, p-1$ and multiply each by $a$:
\[
  a \cdot 1,\ a \cdot 2,\ \dots,\ a \cdot (p-1) \pmod p.
\]
I claim these are the same $p-1$ residues, just shuffled. None is $0$, since $p \nmid a$ and
$p \nmid k$ for $1 \le k \le p-1$, so $p \nmid ak$. And they are distinct: if $ai \equiv aj$ then,
multiplying by $a^{-1}$ (which exists because $\gcd(a,p)=1$), $i \equiv j$. So the list
$\set{a, 2a, \dots, (p-1)a}$ is a permutation of $\set{1, 2, \dots, p-1}$ modulo $p$.

Multiply all of them together. The two products are congruent because the sets match:
\[
  (a \cdot 1)(a \cdot 2)\cdots(a \cdot (p-1)) \equiv 1 \cdot 2 \cdots (p-1) \pmod p,
\]
that is, $a^{p-1}\,(p-1)! \equiv (p-1)! \pmod p$. Since $(p-1)!$ is a product of numbers all coprime
to $p$, it is invertible; cancel it to get $a^{p-1} \equiv 1 \pmod p$.
\end{proof}

\begin{example}[Fermat at $p = 7$]
\label{ch04:ex-flt-7}
Take $a = 3$, $p = 7$. The theorem predicts $3^6 \equiv 1 \pmod 7$. Compute the powers:
$3^1 = 3$, $3^2 = 9 \equiv 2$, $3^3 \equiv 6$, $3^4 \equiv 18 \equiv 4$, $3^5 \equiv 12 \equiv 5$,
$3^6 \equiv 15 \equiv 1 \pmod 7$. It lands on $1$ exactly at the exponent $p - 1 = 6$, and notice
the powers cycled through \emph{all} of $\set{1,2,3,4,5,6}$ on the way --- $3$ is a primitive root
mod $7$, which we return to in Section~\ref{ch04:primitive-roots}.
\end{example}

Euler generalized this from prime moduli to all moduli by replacing $p - 1$ with $\varphi(n)$.

\begin{theorem}[Euler's theorem]
\label{ch04:thm-euler}
If $\gcd(a, n) = 1$, then
\[
  a^{\varphi(n)} \equiv 1 \pmod n.
\]
\end{theorem}

\begin{proof}
The argument is Fermat's, run over the unit group instead of all nonzero residues. Let
$u_1, u_2, \dots, u_{\varphi(n)}$ be the units of $\Zn{n}$ (the residues coprime to $n$). Multiply
each by $a$. Because $a$ is a unit, the map $u \mapsto au$ sends units to units (a product of units
is a unit) and is injective ($au_i \equiv au_j \Rightarrow u_i \equiv u_j$ after multiplying by
$a^{-1}$). An injective map of a finite set to itself is a bijection, so $\set{au_1, \dots,
au_{\varphi(n)}}$ is just the units permuted. Multiply everything:
\[
  \prod_i (a u_i) \equiv \prod_i u_i \pmod n
  \quad\Longrightarrow\quad
  a^{\varphi(n)} \prod_i u_i \equiv \prod_i u_i \pmod n.
\]
The product $\prod_i u_i$ is a unit (units are closed under multiplication), so cancel it:
$a^{\varphi(n)} \equiv 1 \pmod n$.
\end{proof}

Fermat is the special case $n = p$, where $\varphi(p) = p - 1$.

\begin{example}[Euler at $n = 10$]
\label{ch04:ex-euler-10}
Here $\varphi(10) = \varphi(2)\varphi(5) = 1 \cdot 4 = 4$. Take $a = 3$:
$3^4 = 81 \equiv 1 \pmod{10}$. So the last digit of every power of $3$ cycles with period dividing
$4$: $3, 9, 7, 1, 3, 9, 7, 1, \dots$ --- which is why $3^{2026}$ ends in the same digit as
$3^{2026 \bmod 4} = 3^2 = 9$.
\end{example}

\begin{whybox}
This is the engine of RSA correctness. With $n = pq$ and $ed \equiv 1 \pmod{\varphi(n)}$, write
$ed = 1 + k\varphi(n)$. For a message $m$ coprime to $n$,
\[
  m^{ed} = m^{1 + k\varphi(n)} = m \cdot \left(m^{\varphi(n)}\right)^k \equiv m \cdot 1^k = m
  \pmod n
\]
by Euler's theorem. Encrypting with $e$ and decrypting with $d$ are exact inverses. (The leftover
case where $m$ shares a factor with $n$ is handled by the Chinese Remainder Theorem in the next
section --- the identity $m^{ed} \equiv m$ holds for \emph{all} $m$, not just the units, which is a
nice subtlety we will nail down.)
\end{whybox}

\begin{attackbox}
Fermat's theorem is also the basis of the cheapest \emph{primality} screen --- and of a classic
trap. The contrapositive says: if $a^{n-1} \not\equiv 1 \pmod n$ for some $a$ with $\gcd(a,n)=1$,
then $n$ is definitely composite. This is the \emph{Fermat test}. But the converse fails. There are
composite $n$ that satisfy $a^{n-1} \equiv 1 \pmod n$ for \emph{every} $a$ coprime to $n$ --- the
\emph{Carmichael numbers}, the smallest being $561 = 3 \cdot 11 \cdot 17$. A Fermat test will swear
$561$ is prime for every base it cannot detect through a shared factor. Carmichael numbers are
infinite in supply, so the Fermat test alone is not safe for cryptographic prime generation. The
fix is the Miller-Rabin test, which exploits square roots of $1$ to catch Carmichael numbers; we
build it in the chapter on generating primes. The lesson here: a theorem's converse is not a free
gift, and assuming it can be a security hole.
\end{attackbox}

\section{The Chinese Remainder Theorem}
\label{ch04:crt}

Suppose you want a number that leaves remainder $2$ when divided by $3$, remainder $3$ when divided
by $5$, and remainder $2$ when divided by $7$. The Chinese Remainder Theorem (CRT) says there is
exactly one such number modulo $3 \cdot 5 \cdot 7 = 105$, and it hands you a recipe to find it. The
result is ancient --- it appears in Sunzi's \emph{Mathematical Classic}, roughly the 3rd--5th
century CE, in essentially the puzzle just stated.

\begin{theorem}[CRT, two-modulus form]
\label{ch04:thm-crt}
Let $m$ and $n$ be coprime. Then for any targets $a, b$ the system
\[
  x \equiv a \pmod m, \qquad x \equiv b \pmod n
\]
has a unique solution modulo $mn$. The map
\[
  \Zn{mn} \longrightarrow \Zn{m} \times \Zn{n}, \qquad
  x \longmapsto (x \bmod m,\ x \bmod n)
\]
is a bijection that respects addition and multiplication.
\end{theorem}

\begin{proof}
\emph{Existence and the formula.} Since $\gcd(m, n) = 1$, Bezout gives $M m + N n = 1$ for some
integers $M, N$. Note $Nn \equiv 1 \pmod m$ and $\equiv 0 \pmod n$; symmetrically $Mm \equiv 0
\pmod m$ and $\equiv 1 \pmod n$. Set
\[
  x = a \cdot (Nn) + b \cdot (Mm).
\]
Modulo $m$ the second term vanishes and $Nn \equiv 1$, so $x \equiv a$. Modulo $n$ the first term
vanishes and $Mm \equiv 1$, so $x \equiv b$. Done.

\emph{Uniqueness.} If $x$ and $x'$ both solve the system, then $m \mid (x - x')$ and
$n \mid (x - x')$. Since $\gcd(m, n) = 1$, their product divides the difference (any common multiple
of coprime numbers is a multiple of their product), so $mn \mid (x - x')$, i.e. $x \equiv x' \pmod
{mn}$.

The map is clearly a ring homomorphism (reduction mod $m$ and mod $n$ both respect $+$ and
$\times$), it is injective by uniqueness, and both sides have $mn$ elements, so it is a bijection.
\end{proof}

\begin{example}[Sunzi's puzzle, worked]
\label{ch04:ex-crt-sunzi}
Solve $x \equiv 2 \pmod 3$, $x \equiv 3 \pmod 5$, $x \equiv 2 \pmod 7$. Do it pair by pair.

First combine mod $3$ and mod $5$ into mod $15$. We need $x \equiv 2 \pmod 3$ and $x \equiv 3
\pmod 5$. Bezout for $(3,5)$: $2 \cdot 3 - 1 \cdot 5 = 1$, so $-1 \cdot 5 \equiv 1 \pmod 3$ and
$2 \cdot 3 \equiv 1 \pmod 5$. Then
\[
  x = 2 \cdot (-5) + 3 \cdot (6) = -10 + 18 = 8, \qquad x \equiv 8 \pmod{15}.
\]
Check: $8 = 2 \cdot 3 + 2$ and $8 = 1 \cdot 5 + 3$. Good.

Now combine $x \equiv 8 \pmod{15}$ with $x \equiv 2 \pmod 7$. Bezout for $(15, 7)$:
$1 = 15 - 2 \cdot 7$, so $15 \equiv 1 \pmod 7$ and $-2 \cdot 7 \equiv 1 \pmod{15}$. Then
\[
  x = 8 \cdot (-14) + 2 \cdot (15) = -112 + 30 = -82 \equiv -82 + 105 = 23 \pmod{105}.
\]
So $x \equiv 23 \pmod{105}$. Verify all three: $23 = 7\cdot3 + 2$, $23 = 4 \cdot 5 + 3$,
$23 = 3 \cdot 7 + 2$. All correct.
\end{example}

\begin{intuitionbox}
The CRT says a residue modulo $mn$ is nothing more than a \emph{pair} of residues, one mod $m$ and
one mod $n$ --- and you can add or multiply the pairs coordinatewise. A single odometer counting to
$mn$ is indistinguishable from two independent odometers of lengths $m$ and $n$ spinning in
parallel, \emph{provided their lengths are coprime} so they never resynchronize early. Coprimality
is what guarantees all $mn$ pairs actually occur.
\end{intuitionbox}

\begin{figure}[H]
\centering
\begin{tikzpicture}[font=\small,>=stealth,scale=0.95]
  % Unified ring box
  \node[draw=structurepurple!80,fill=structurepurple!12,thick,rounded corners=4pt,minimum width=2.6cm,minimum height=1.3cm,align=center] (big) at (0,0) {
    \textbf{Unified Ring}\\
    $x \in \Zn{6}$
  };

  % Independent factor rings
  \node[draw=sourceblue!80,fill=sourceblue!12,thick,rounded corners=4pt,minimum width=2.8cm,minimum height=1.1cm,align=center] (a) at (6,0.9) {
    \textbf{Factor 1:}\\
    $x_1 \equiv x \pmod 2 \in \Zn{2}$
  };

  \node[draw=targetgreen!80,fill=targetgreen!12,thick,rounded corners=4pt,minimum width=2.8cm,minimum height=1.1cm,align=center] (b) at (6,-0.9) {
    \textbf{Factor 2:}\\
    $x_2 \equiv x \pmod 3 \in \Zn{3}$
  };

  % Arrows representing isomorphism
  \draw[->,very thick,sourceblue!90!black] (big.east) to[out=25,in=180] node[above,font=\tiny\bfseries]{$\bmod 2$} (a.west);
  \draw[->,very thick,targetgreen!90!black] (big.east) to[out=-25,in=180] node[below,font=\tiny\bfseries]{$\bmod 3$} (b.west);

  \draw[<-,thick,dashed,quotientorange] (big.south east) to[out=-40,in=200] node[below,font=\tiny\bfseries]{$\text{CRT}^{-1}: \text{Garner / B\'ezout}$} (b.south west);

  % Table / worked example below
  \node[draw=gray!40,fill=gray!5,rounded corners=3pt,thick,align=center,font=\footnotesize,inner sep=5pt] at (3.0,-2.3) {
    \textbf{Ring Isomorphism:} $\Zn{6} \cong \Zn{2} \times \Zn{3}$ since $\gcd(2,3)=1$.\\
    $5 \leftrightarrow (1, 2), \quad 4 \leftrightarrow (0, 1), \quad 5\cdot 4 = 20 \equiv 2 \leftrightarrow (1\cdot 0, 2\cdot 1) = (0, 2) \equiv 2$.
  };
\end{tikzpicture}
\caption{The Chinese Remainder Theorem as a ring isomorphism. Because 2 and 3 are coprime, computing modulo 6 is structurally identical to running two independent, parallel computations mod 2 and mod 3. Operations can be evaluated in the smaller, faster sub-rings and recombined at the end.}
\label{ch04:fig-crt-bijection}
\end{figure}

\subsection*{Totient is multiplicative}

\begin{corollary}
\label{ch04:cor-totient-mult}
If $\gcd(m, n) = 1$ then $\varphi(mn) = \varphi(m)\,\varphi(n)$.
\end{corollary}

\begin{proof}
The CRT bijection $\Zn{mn} \to \Zn{m} \times \Zn{n}$ respects multiplication, so it carries units to
units: $x$ is invertible mod $mn$ iff both coordinates are invertible. Hence it restricts to a
bijection between $\Zns{mn}$ and $\Zns{m} \times \Zns{n}$. Counting both sides,
$\varphi(mn) = \varphi(m)\varphi(n)$.
\end{proof}

This is the missing piece that completes the totient formula of Section~\ref{ch04:congruences}: a
general $n$ splits into coprime prime powers, and $\varphi$ multiplies across them.

\begin{applicationbox}
CRT makes RSA decryption about four times faster, and every serious RSA implementation uses it. The
private key secretly stores $p$ and $q$. To compute $m = c^d \bmod n$, instead of one exponentiation
modulo the $2k$-bit number $n$, you do two exponentiations modulo the $k$-bit numbers $p$ and $q$:
\[
  m_p = c^{d \bmod (p-1)} \bmod p, \qquad m_q = c^{d \bmod (q-1)} \bmod q,
\]
then glue $m_p, m_q$ back together with the CRT formula to recover $m \bmod n$. Each half-size
exponentiation is roughly four times cheaper (modular multiplication scales quadratically in the
operand size), and the gluing is negligible. The catch: a fault during one of the two halves leaks
the factorization. If $m_p$ is computed correctly but $m_q$ is corrupted, then the faulty signature
$s'$ satisfies $\gcd(s'^e - m,\ n) = p$ exactly, handing an attacker the secret prime. This
\emph{Bellcore fault attack} is why CRT-RSA implementations verify the result before releasing it.
\end{applicationbox}

\begin{implementbox}
A general CRT solver over many coprime moduli, building the answer one congruence at a time.

\begin{lstlisting}[language=Python]
def crt(residues, moduli):
    # solve x = residues[i] (mod moduli[i]); moduli pairwise coprime
    x, M = 0, 1
    for a, m in zip(residues, moduli):
        # current x solves the first batch mod M; fold in (a mod m)
        g, p, q = egcd(M, m)        # M*p + m*q = 1  (g == 1 required)
        if g != 1:
            raise ValueError("moduli not coprime")
        x = (a * M * p + x * m * q) % (M * m)
        M *= m
    return x % M

print(crt([2, 3, 2], [3, 5, 7]))   # -> 23
\end{lstlisting}

\textbf{Learning only.} No constant-time guarantees and no fault checking. Real CRT-RSA code
recombines $m_p$ and $m_q$ with the precomputed inverse $q^{-1} \bmod p$ (Garner's formula) and
then re-encrypts to verify, precisely to block the Bellcore fault attack above.
\end{implementbox}

\section{Primitive Roots}
\label{ch04:primitive-roots}

Euler's theorem tells us $a^{\varphi(n)} \equiv 1$, but $\varphi(n)$ might not be the \emph{first}
exponent that returns to $1$. The first one has a name.

\begin{definition}[Multiplicative order]
\label{ch04:def-order}
For a unit $a$ modulo $n$, the \emph{order} $\ord_n(a)$ is the smallest positive integer $k$ with
$a^k \equiv 1 \pmod n$.
\end{definition}

\begin{lemma}
\label{ch04:lem-order-divides}
$a^m \equiv 1 \pmod n$ if and only if $\ord_n(a) \mid m$. In particular $\ord_n(a) \mid \varphi(n)$.
\end{lemma}

\begin{proof}
Let $d = \ord_n(a)$. Divide $m = dq + r$ with $0 \le r < d$. Then
$a^m = (a^d)^q a^r \equiv a^r \pmod n$, so $a^m \equiv 1$ iff $a^r \equiv 1$. By minimality of $d$,
$0 \le r < d$ forces $r = 0$, i.e. $d \mid m$. The ``in particular'' follows from Euler's theorem,
which gives $a^{\varphi(n)} \equiv 1$.
\end{proof}

\begin{definition}[Primitive root]
\label{ch04:def-primroot}
A unit $g$ modulo $n$ is a \emph{primitive root} if $\ord_n(g) = \varphi(n)$ --- the largest order
possible. Then the powers $g^0, g^1, \dots, g^{\varphi(n)-1}$ run through \emph{every} unit exactly
once: $g$ \emph{generates} $\Zns{n}$, and the group is \emph{cyclic}.
\end{definition}

\begin{example}[Primitive root mod $7$]
\label{ch04:ex-primroot-7}
We saw in Example~\ref{ch04:ex-flt-7} that the powers of $3$ mod $7$ are
$3, 2, 6, 4, 5, 1$ --- all six units before returning to $1$. So $\ord_7(3) = 6 = \varphi(7)$ and
$3$ is a primitive root. Compare $2$: its powers are $2, 4, 1, 2, 4, 1, \dots$, order $3$, a proper
divisor of $6$. So $2$ is not a primitive root mod $7$. The orders of the six units are
$\ord(1)=1$, $\ord(6)=2$, $\ord(2)=\ord(4)=3$, $\ord(3)=\ord(5)=6$ --- and every order divides $6$,
as Lemma~\ref{ch04:lem-order-divides} promised.
\end{example}

Not every modulus has a primitive root. The ones that do are exactly $n \in \set{1, 2, 4, p^k,
2p^k}$ for an odd prime $p$ --- a theorem of Gauss we will not prove here. The case we care about
most is clean.

\begin{theorem}[Primitive roots mod a prime]
\label{ch04:thm-primroot-prime}
For every prime $p$, the group $\Zns{p}$ is cyclic: a primitive root exists, and there are exactly
$\varphi(p-1)$ of them.
\end{theorem}

The proof uses the fact that a degree-$d$ polynomial over $\Zn{p}$ has at most $d$ roots (a field
property we develop in the finite-fields chapter), so we defer the full argument. The structural
payoff is enormous: once a generator $g$ is fixed, every nonzero residue is $g^k$ for a unique
exponent $k \in \set{0, \dots, p-2}$, and that exponent is the \emph{discrete logarithm}
$\dlog_g(x)$.

\begin{intuitionbox}
A primitive root turns the multiplicative group $\Zns{p}$ into a clock with $p-1$ positions.
Multiplying two units corresponds to \emph{adding} their discrete logs and wrapping around. So
multiplication mod $p$ is secretly just addition mod $p-1$ in disguise --- once you know the
dictionary $x \leftrightarrow \dlog_g(x)$. The catch, and the entire foundation of Diffie-Hellman,
is that going from $x$ to $g^x$ is fast while going back from $g^x$ to $x$ (computing the discrete
log) appears to be exponentially hard for well-chosen $p$.
\end{intuitionbox}

\begin{figure}[H]
\centering
\begin{tikzpicture}[font=\small,>=stealth,scale=1.05]
  \def\R{1.8}
  % 6 positions: powers of 3 mod 7: 3^0=1, 3^1=3, 3^2=2, 3^3=6, 3^4=4, 3^5=5
  \foreach \i/\val/\exp in {0/1/0, 1/3/1, 2/2/2, 3/6/3, 4/4/4, 5/5/5} {
    \pgfmathsetmacro\ang{90 - \i*60}
    \node[circle,draw=structurepurple!80,fill=structurepurple!12,thick,minimum size=24pt,align=center]
      (n\i) at (\ang:\R) {\textbf{\val}};
    \node[font=\tiny,gray] at (\ang:\R+0.42) {$3^{\exp}$};
  }

  \foreach \i in {0,1,2,3,4,5} {
    \pgfmathsetmacro\j{int(mod(\i+1,6))}
    \draw[->,very thick,structurepurple!90!black] (n\i) to[bend left=20] 
      node[midway,fill=white,inner sep=1pt,font=\tiny\bfseries,text=structurepurple!90!black]{$\times 3$} (n\j);
  }

  \node[align=center,font=\footnotesize\bfseries,structurepurple] at (0,0.2) {Cyclic Group\\$\Zns{7} = \langle 3 \rangle$};
  \node[align=center,font=\tiny,gray] at (0,-0.4) {$\ord_7(3) = 6 = \varphi(7)$};

  \node[draw=structurepurple!40,fill=structurepurple!5,rounded corners=3pt,font=\tiny,align=center,inner sep=4pt] at (0,-2.6) {
    Each step advances the exponent by $+1 \pmod 6$.\\
    Multiplication in the group $\Zns{7}$ is equivalent to addition on the exponent clock $\Zn{6}$.
  };
\end{tikzpicture}
\caption{The cyclic orbit of the primitive root $g = 3$ modulo 7. The powers $3^k \bmod 7$ cycle through every single non-zero residue before returning to 1, providing the algebraic foundation for the Diffie-Hellman discrete logarithm problem.}
\label{ch04:fig-primroot-clock}
\end{figure}

\begin{example}[Counting primitive roots mod $11$]
\label{ch04:ex-primroot-11}
Here $\varphi(11) = 10$, so the count of primitive roots is $\varphi(10) = \varphi(2)\varphi(5)
= 1 \cdot 4 = 4$. Test $g = 2$: its order must divide $10$, so it is $1, 2, 5,$ or $10$. Check the
proper divisors --- $2^2 = 4 \neq 1$, $2^5 = 32 \equiv 10 \neq 1 \pmod{11}$ --- neither hits $1$,
so $\ord_{11}(2) = 10$ and $2$ is a primitive root. The four primitive roots are the powers $2^k$
with $\gcd(k, 10) = 1$, namely $k \in \set{1, 3, 7, 9}$: $2^1 = 2$, $2^3 = 8$, $2^7 \equiv 7$,
$2^9 \equiv 6 \pmod{11}$. So $\set{2, 6, 7, 8}$.
\end{example}

\begin{connectionbox}
Why test only the divisors of $\varphi(n)$, not all exponents up to $\varphi(n)$? Because
Lemma~\ref{ch04:lem-order-divides} pins the order to a divisor of $\varphi(n)$. To confirm $g$ is a
primitive root it suffices to check $g^{\varphi(n)/q} \not\equiv 1$ for each prime $q$ dividing
$\varphi(n)$ --- if none of those drop to $1$, no proper divisor can be the order. This is the
standard fast primitive-root test, and it needs the factorization of $\varphi(n)$. The probability
side of the story (a random residue is a primitive root with probability $\varphi(p-1)/(p-1)$,
which the prime-number theorem keeps bounded away from zero) lives in the companion probability
book.
\end{connectionbox}

\begin{applicationbox}
Diffie-Hellman key exchange lives in exactly this cyclic structure. Alice and Bob publicly agree on
a large prime $p$ and a generator $g$ of (a large prime-order subgroup of) $\Zns{p}$. Alice picks a
secret $a$ and sends $g^a \bmod p$; Bob picks secret $b$ and sends $g^b \bmod p$. Each raises the
other's value to their own secret, and both arrive at $g^{ab} \bmod p$ --- the shared key. Eve sees
$g$, $g^a$, $g^b$ but to find $g^{ab}$ she would seemingly need a discrete logarithm, which is
infeasible for a well-chosen $p$ of a couple thousand bits. The whole protocol is the observation
from the intuition box --- exponentiation is easy, its inverse is hard --- dressed up as a
conversation. Elliptic-curve Diffie-Hellman replaces $\Zns{p}$ with the points on a curve but reuses
this identical cyclic-group skeleton.
\end{applicationbox}

\section{Putting It Together: A Tiny RSA by Hand}
\label{ch04:rsa-toy}

Every tool in this chapter now assembles into one worked example. We build a complete RSA keypair
with $p = 61$, $q = 53$ and run a message through it (Figure~\ref{ch04:fig-rsa-roundtrip}),
computing each piece with the methods above.

\textbf{Key generation.}
\begin{enumerate}
\item Modulus: $n = pq = 61 \cdot 53 = 3233$.
\item Totient: $\varphi(n) = (p-1)(q-1) = 60 \cdot 52 = 3120$
      (Example~\ref{ch04:ex-totient-rsa}).
\item Public exponent: choose $e = 17$. It is coprime to $3120$: since
      $3120 = 2^4 \cdot 3 \cdot 5 \cdot 13$ has no factor of $17$, $\gcd(17, 3120) = 1$, confirmed
      by extended Euclid below.
\item Private exponent: $d = e^{-1} \bmod \varphi(n) = 17^{-1} \bmod 3120$. Run extended Euclid:
      \[
      \begin{aligned}
      3120 &= 183 \cdot 17 + 9 \\
        17 &= 1 \cdot 9 + 8 \\
         9 &= 1 \cdot 8 + 1.
      \end{aligned}
      \]
      Back-substitute: $1 = 9 - 8 = 9 - (17 - 9) = 2\cdot 9 - 17 = 2(3120 - 183\cdot 17) - 17
      = 2\cdot 3120 - 367\cdot 17$. So $-367 \cdot 17 \equiv 1 \pmod{3120}$, giving
      $d \equiv -367 \equiv 2753 \pmod{3120}$.
      Check: $17 \cdot 2753 = 46801 = 15 \cdot 3120 + 1$. Correct.
\end{enumerate}
The public key is $(n, e) = (3233, 17)$; the private key is $d = 2753$ (with $p, q$ kept for the
CRT speedup).

\textbf{Encrypting.} Let the message be $m = 65$ (say, the letter \texttt{A} in ASCII). The
ciphertext is
\[
  c = m^e \bmod n = 65^{17} \bmod 3233.
\]
Compute by square-and-multiply, reducing mod $3233$ at every step. Powers of $65$:
$65^2 = 4225 \equiv 992$; $65^4 \equiv 992^2 = 984064 \equiv 984064 - 304\cdot 3233 = 1232$;
$65^8 \equiv 1232^2 = 1517824 \equiv 1517824 - 469\cdot 3233 = 1547$;
$65^{16} \equiv 1547^2 = 2393209 \equiv 2393209 - 740\cdot 3233 = 789$. Since
$17 = 16 + 1$,
\[
  c = 65^{16} \cdot 65 \equiv 789 \cdot 65 = 51285 \equiv 51285 - 15\cdot 3233 = 2790 \pmod{3233}.
\]
So $c = 2790$.

\textbf{Decrypting with the CRT speedup.} We must recover $m = c^d \bmod n = 2790^{2753} \bmod 3233$.
Done directly this is a $2753$-exponent slog, so use the CRT route --- exactly what real RSA does.
Reduce the private exponent modulo each factor's totient:
\[
  d_p = d \bmod (p-1) = 2753 \bmod 60 = 53, \qquad
  d_q = d \bmod (q-1) = 2753 \bmod 52 = 49.
\]
Reduce the ciphertext into each prime: $2790 \bmod 61 = 45$ and $2790 \bmod 53 = 34$. Then by Fermat
the two half-size exponentiations are
\[
  m_p = 45^{53} \bmod 61 = 4, \qquad m_q = 34^{49} \bmod 53 = 12.
\]
Glue them with Garner's form of the CRT. The formula needs the inverse \emph{of $p$ modulo $q$}:
\[
  m = m_p + p\cdot\big[(m_q - m_p)\,(p^{-1}\bmod q)\bmod q\big].
\]
A tempting mistake is to grab the wrong inverse. Garner wants $p^{-1} \bmod q$, i.e. $61^{-1} \bmod
53$; since $61 \equiv 8 \pmod{53}$ this is $8^{-1} \bmod 53$. If you carelessly used some other
value such as $12$, you would get $(12-4)\cdot 12 = 96 \equiv 43 \pmod{53}$ and
$m = 4 + 61\cdot 43 = 2627$ --- a legal-looking residue (it is below $n = 3233$) that is silently
\emph{wrong}, because the round trip must return $m = 65$. The genuine inverse is $8^{-1} \equiv 20
\pmod{53}$ (check $8\cdot 20 = 160 = 3\cdot 53 + 1$). With it,
\[
  m = 4 + 61\cdot\big[(12 - 4)\cdot 20 \bmod 53\big]
    = 4 + 61\cdot\big[160 \bmod 53\big] = 4 + 61\cdot 1 = 65.
\]
The round trip closes, and Euler's theorem guaranteed it would.

\begin{figure}[H]
\centering
\begin{tikzpicture}[scale=1, every node/.style={font=\footnotesize}]
  \node[draw=aliceblue, fill=aliceblue!8, rounded corners, minimum width=2.4cm,
        minimum height=1.2cm, align=center] (alice) at (0,0)
        {\textbf{Alice}\\$m = 65$};
  \node[draw=bobgreen, fill=bobgreen!8, rounded corners, minimum width=2.4cm,
        minimum height=1.2cm, align=center] (bob) at (8,0)
        {\textbf{Bob}\\holds $d = 2753$};
  % encryption arrow
  \draw[->,>=Stealth,thick] (alice.east) -- node[above,align=center]
        {$c = m^{e}\bmod n$\\$= 65^{17}\bmod 3233 = 2790$} (bob.west);
  % public key label
  \node[below=0.15cm of alice, align=center, structurepurple]
        {uses public key\\$(n,e)=(3233,17)$};
  % decryption label
  \node[below=0.15cm of bob, align=center, structurepurple]
        {$m = c^{d}\bmod n$\\$= 2790^{2753}\bmod 3233 = 65$};
  % Eve
  \node[draw=evered, fill=evered!8, rounded corners, minimum width=2.2cm,
        minimum height=0.9cm, align=center] (eve) at (4,-2.6)
        {\textbf{Eve}\\sees $c=2790$, $(n,e)$};
  \draw[->,>=Stealth,evered,dashed] (4,-0.5) -- (eve.north);
  \node[evered, right=0.1cm of eve, align=left]
        {needs to factor $n=pq$\\to learn $d$ --- infeasible at scale};
\end{tikzpicture}
\caption{The toy RSA round trip. Encryption with the public exponent and decryption with the
private exponent are exact inverses (Euler's theorem); Eve, lacking the factorization of $n$,
cannot reconstruct $d$.}
\label{ch04:fig-rsa-roundtrip}
\end{figure}

\begin{walkbox}
The lesson of that near-miss: Garner recombination is exactly where careless CRT-RSA code goes
wrong, and the bug is silent --- you get \emph{a} residue, just not the right one. The fix is to use
the genuine inverse of $p \bmod q$. Work the whole decryption yourself: recompute
$45^{53}\bmod 61$ and $34^{49}\bmod 53$ by square-and-multiply (the chains are short, five or six
squarings each), confirm $m_p = 4$, $m_q = 12$, and verify Garner with the coefficient $20$ lands on
$65$. If it does, you have personally executed every theorem in this chapter --- division with
remainder, extended Euclid for $d$, Fermat for the subgroup exponents, and CRT for the recombination
--- on real numbers. That is the whole toolkit working as one machine.
\end{walkbox}

\section{Exercises}
\label{ch04:exercises}

\begin{exercise}[Warm-up gcd and Bezout]
Compute $\gcd(252, 198)$ with the Euclidean algorithm, showing each division. Then run extended
Euclid to find integers $x, y$ with $252x + 198y = \gcd(252,198)$, and verify your answer by direct
substitution.
\end{exercise}

\begin{exercise}[Modular inverse]
Compute $11^{-1} \bmod 26$ (the inverse that ``undoes'' an affine cipher with multiplier $11$ on the
26-letter alphabet). Show the extended-Euclid steps and check that the product is $\equiv 1$.
Then explain why $13^{-1} \bmod 26$ does \emph{not} exist.
\end{exercise}

\begin{exercise}[Fermat in action]
Use Fermat's little theorem to compute $7^{222} \bmod 11$ without a calculator, by first reducing
the exponent modulo $10$. Then compute $2^{1000} \bmod 13$ the same way.
\end{exercise}

\begin{exercise}[Totient practice]
Compute $\varphi(360)$, $\varphi(1000)$, and $\varphi(p^2 q)$ for distinct primes $p, q$ (give a
formula for the last). For which $n \le 30$ is $\varphi(n)$ odd? Prove your characterization.
\end{exercise}

\begin{exercise}[CRT by hand]
Find the smallest positive integer $x$ with $x \equiv 1 \pmod 4$, $x \equiv 2 \pmod 5$, and
$x \equiv 3 \pmod 7$. Show the pairwise combining steps, and state the full set of solutions modulo
$140$.
\end{exercise}

\begin{exercise}[Order and primitive roots]
For the prime $p = 13$: (a) compute the order of $2$, $3$, and $5$ modulo $13$; (b) determine which
of them are primitive roots; (c) using a primitive root you found, list the discrete logarithm of
every nonzero residue. How many primitive roots does $13$ have, and does your count match
$\varphi(12)$?
\end{exercise}

\begin{exercise}[Break a toy RSA]
An adversary intercepts the public key $(n, e) = (3233, 17)$ and, by factoring the small modulus,
recovers $p = 61$, $q = 53$. Reconstruct the private exponent $d$ from scratch with extended Euclid,
then decrypt the ciphertext $c = 2790$ to recover the plaintext. (You have just performed the entire
RSA break that large primes are designed to prevent.)
\end{exercise}

\begin{exercise}[Recover the factors from the totient]
You are told that $n = 8051$ is a product of two primes and that $\varphi(n) = 7872$. Without
trial division, recover $p$ and $q$ by setting up and solving a quadratic in $p + q$ and $pq$, as in
Example~\ref{ch04:ex-totient-rsa}. Explain in one sentence why this shows that computing
$\varphi(n)$ is as hard as factoring.
\end{exercise}

\begin{exercise}[Carmichael number]
Verify that $561 = 3 \cdot 11 \cdot 17$ is a Fermat pseudoprime to base $2$ --- that is,
$2^{560} \equiv 1 \pmod{561}$ --- by computing $2^{560}$ modulo $3$, $11$, and $17$ separately
(Fermat makes each trivial) and recombining with CRT. Why does this make $561$ dangerous to a naive
primality test?
\end{exercise}

\begin{exercise}[Euclid's lemma, used in anger]
Prove that $\sqrt{2}$ is irrational using the Fundamental Theorem of Arithmetic: suppose
$\sqrt 2 = a/b$ in lowest terms, square it, and derive a contradiction from the parity of the prime
factorization of $2b^2 = a^2$. Identify exactly where Euclid's lemma (or unique factorization) is
the load-bearing step.
\end{exercise}

\begin{exercise}[Prove a divisibility identity]
Prove that for all integers $n \ge 1$, $6 \mid n^3 - n$. (Hint: factor $n^3 - n = (n-1)n(n+1)$ and
argue about three consecutive integers, or use Fermat's little theorem at the primes $2$ and $3$
together with CRT.)
\end{exercise}

\begin{exercise}[Finish the worked decryption]
Complete the CRT recombination from Section~\ref{ch04:rsa-toy} cleanly. Recompute
$m_p = 45^{53} \bmod 61$ and $m_q = 34^{49} \bmod 53$ with square-and-multiply, then apply Garner's
formula $m = m_p + p\cdot\big[(m_q - m_p)\,(p^{-1}\bmod q)\bmod q\big]$ with the correct coefficient
$p^{-1} \equiv 20 \pmod{53}$, and confirm you recover $m = 65$. Explain precisely why the value
$2627$ appeared in the text and what single number was wrong.
\end{exercise}
